\chapter{TỔNG QUAN TÀI LIỆU}
\label{chap:related}

\section{Phân tích cảm xúc là gì}

Phân tích cảm xúc, còn được gọi là khai thác ý kiến, là nghiên cứu tính toán về ý kiến, cảm xúc, tâm trạng và thái độ của mọi người được thể hiện trong văn bản \cite{liu2012sentiment}.

\section{Các phương pháp truyền thống}

Các cách tiếp cận sớm cho phân tích cảm xúc dựa vào:
\begin{itemize}
    \item \textbf{Phương pháp dựa trên từ điển}: Sử dụng các từ điển cảm xúc với điểm số cảm xúc được gán sẵn
    \item \textbf{Phương pháp học máy}: Đào tạo các bộ phân loại trên các đặc trưng được tạo thủ công
\end{itemize}

\subsection{Phương pháp dựa trên từ điển}

Phương pháp dựa trên từ điển sử dụng các từ điển cảm xúc như SenticNet \cite{cambria2016senticnet} và VADER \cite{hutto2014vader} để xác định cảm xúc của văn bản bằng cách tổng hợp điểm số cảm xúc của các từ và cụm từ.

\subsection{Phương pháp học máy truyền thống}

Các phương pháp học máy truyền thống bao gồm Naive Bayes, Support Vector Machines (SVM), và Random Forest đã được sử dụng rộng rãi cho phân loại cảm xúc \cite{pang2002thumbs}.

\section{Phương pháp Deep Learning}

\subsection{Word Embeddings}

Word embeddings như Word2Vec \cite{mikolov2013distributed} và GloVe \cite{pennington2014glove} đã cải thiện đáng kể hiệu suất phân tích cảm xúc bằng cách nắm bắt các mối quan hệ ngữ nghĩa giữa các từ.

\subsection{Recurrent Neural Networks}

RNNs, LSTMs, và GRUs đã được áp dụng thành công cho các tác vụ phân tích cảm xúc \cite{zhou2016text}.

\subsection{Transformer Models}

Các mô hình dựa trên transformer gần đây như BERT \cite{devlin2018bert}, RoBERTa \cite{liu2019roberta} và VNLPT \cite{van2022vnlp} đã đạt được kết quả tiên tiến trong phân loại cảm xúc.

\section{Phân tích cảm xúc tiếng Việt}

Nghiên cứu về phân tích cảm xúc tiếng Việt đã phát triển cùng với sự phát triển của các mô hình ngôn ngữ và bộ dữ liệu tiếng Việt \cite{van2022vnlp}.
